% DEFINED:   sec:stability, sec:gradcheck, sec:autograd-note
% RESOLVED:  sec:logits-logprobs, sec:mlp-backprop
% FORWARD:   (none)

\section{Numerical stability is the gauge freedom, reused}\label{sec:stability}

$e^z$ overflows for moderately large $z$. The fix for softmax: subtract
the largest logit first,
\[
  \softmax(\vz) = \softmax\!\left(\vz - \max_i z_i\right).
\]
This is \emph{exactly} the additive-constant freedom from
\Cref{sec:logits-logprobs}. Softmax does not care about a shared shift
of the logits, so we spend that freedom to keep the exponentials in a
safe range. Stability here is not a new idea bolted on. It is the gauge
symmetry put to work.

Cross-entropy is computed as $\mathrm{logsumexp}(\vz) - z_y$ rather than
$-\log \mathrm{softmax}(\vz)_y$, which would risk $\log 0$. The sigmoid
is evaluated with a sign branch; binary cross-entropy is computed straight
from the logit, both to avoid overflow and $\log 0$. These are the
primitives the loss classes call:

\begin{lstlisting}
def sigmoid(z):
    if z >= 0.0:
        return 1.0 / (1.0 + math.exp(-z))     # safe: exp(-z), z >= 0
    ez = math.exp(z)                          # safe: exp(z),  z < 0
    return ez / (1.0 + ez)

def logsumexp(zs):
    m = max(zs)
    return m + math.log(sum(math.exp(z - m) for z in zs))

def softmax(zs):
    m = max(zs)
    exps = [math.exp(z - m) for z in zs]
    total = sum(exps)
    return [e / total for e in exps]
\end{lstlisting}

In every case the exponent argument is held at or below $0$, so $e^z$
stays in $[0, 1]$ and never overflows. The max-subtraction in
\texttt{softmax} and \texttt{logsumexp} is the \Cref{sec:logits-logprobs}
gauge freedom; the sign branch in \texttt{sigmoid} is the same trick for
the two-class case.

\section{Trust, but verify: numerical gradients}\label{sec:gradcheck}

Hand-derived gradients are easy to get subtly wrong. There is a slow but
foolproof check: the definition of a derivative itself. For any parameter
$\theta$,
\[
  \frac{\partial L}{\partial \theta}
  \approx
  \frac{L(\theta + \varepsilon) - L(\theta - \varepsilon)}{2\varepsilon}.
\]
\texttt{gradient\_check} computes this central difference for every
parameter and compares it to the analytical gradient from
\texttt{backward}. It perturbs each parameter in place, twice, and reads
off the loss:

\begin{lstlisting}
def gradient_check(net, x, y, eps=1e-5):
    net.zero_grad()
    net.forward(x)
    net.backward(y)
    worst = 0.0
    for values, grads in net.parameters():
        for k in range(len(values)):
            original = values[k]
            values[k] = original + eps
            loss_plus = net.loss_value(x, y)
            values[k] = original - eps
            loss_minus = net.loss_value(x, y)
            values[k] = original
            numerical = (loss_plus - loss_minus) / (2.0 * eps)
            analytical = grads[k]
            denom = max(abs(numerical) + abs(analytical), 1e-12)
            worst = max(worst, abs(numerical - analytical) / denom)
    return worst
\end{lstlisting}

This is too slow to train with: one forward pass per parameter. But it is
the ground truth that keeps the fast method honest. It works for
\emph{any} layer and loss because it only ever touches
\texttt{parameters()}, \texttt{forward}, and \texttt{backward}. Running
it on the four configurations from \Cref{sec:mlp-backprop} (seed~42):

\begin{itemize}
\item Logistic regression (SigmoidBCE): worst relative error $\approx 6 \times 10^{-12}$
\item Softmax regression (SoftmaxCE): $\approx 1 \times 10^{-10}$
\item Tanh MLP (SoftmaxCE): $\approx 7 \times 10^{-9}$
\item ReLU MLP (SigmoidBCE): $\approx 1 \times 10^{-9}$
\end{itemize}

All residuals are near machine precision. One caveat: \texttt{ReLU} has a
kink at $0$ and no derivative there. A finite difference that straddles
the kink will disagree with the analytical sub-gradient. With random
initial weights the probability of a pre-activation landing exactly on
the kink is effectively zero, so the check passes in practice. The
tolerance is anchored on a \texttt{Tanh} network, which is smooth
everywhere; the $10^{-9}$ figure for the ReLU network is real, not a
relaxed threshold masking a bug.

\section{Closing note: from per-layer to per-operation}\label{sec:autograd-note}

Backpropagation here is per-\emph{layer}: each layer carries a local
\texttt{backward}. But nothing forced the unit to be a layer. If every
scalar \emph{operation} (every $+$, every $\times$, every $\tanh$)
carried its own local backward, the same reverse pass would compute
gradients through an arbitrary expression, with no hand-derived layer
gradients at all. That is \textbf{automatic differentiation}, and it is
what every modern framework's autograd engine does
(\cite{baydin2018autodiff}). It is the same idea as this library, at a
finer grain. Building one is the natural next step, and Chapter~2
returns to this template when framing its own natural extensions.
